\section{Conclusion and Future Work}
\label{sec:conclusion}

\subsection{Summary}
\label{sec:conclusion:summary}

We presented a prototype end-to-end system for cross-language scene
text replacement in video: a modular five-stage pipeline in which
two compact networks we designed and trained from scratch (the
Blur Prediction Network and the Alignment Refiner) sit alongside
five pretrained components accessed through stable interfaces.
Evaluation on real video examples shows the feasibility of the
approach as a strong baseline in the glyph-controlled paradigm,
complementary to end-to-end generative video editors that
currently offer superior style preservation at the cost of glyph
accuracy.

\subsection{Model involvement}
\label{sec:conclusion:involvement}

Table~\ref{tab:involvement} summarizes the provenance of each model
in the pipeline, and the nature of our own involvement with it.

\begin{table}[H]
  \centering
  \caption{Models in the pipeline and our involvement with each.}
  \label{tab:involvement}
  \small
  \setlength{\tabcolsep}{4pt}
  \begin{tabular}{lll}
    \toprule
    Model & Stage & Our involvement \\
    \midrule
    PaddleOCR (PP-OCRv4) & S1 detection       & Pre-trained, used as is \\
    CoTracker3           & S1 tracking        & Pre-trained, used as is \\
    AnyText2             & S3 text editing    & Pre-trained, used as is (Gradio server) \\
    SRNet                & S3/S4/S5 inpainting & Pre-trained, used as is \\
    Hi-SAM               & S3/S4/S5 inpainting & Pre-trained, used as is (default backend) \\
    LCM                  & S4 lighting        & Reimplementation of STRIVE's algorithm \\
    \textbf{BPN}         & S4 blur prediction & \textbf{Trained from scratch} on S2-aligned data \\
    \textbf{Alignment Refiner} & S2 drift correction & \textbf{Designed and trained from scratch} \\
    \bottomrule
  \end{tabular}
\end{table}

\subsection{Future work}
\label{sec:conclusion:future}

Several directions follow naturally from the design decisions and
limitations described above.

\begin{itemize}[leftmargin=*]
  \item \textbf{Stronger detection filtering.} The S1 detection
    stage still lets through two kinds of unwanted detections that
    the current wordfreq-based Zipf filter does not catch:
    short real words recognized on textured or geometric
    backgrounds (false positives that happen to look like common
    English), and genuine but unwanted on-screen text such as
    watermarks, timestamps, and tiny background signage. Both
    cause the downstream pipeline to spend effort on regions that
    should not be translated, and in the worst case introduce
    visible artifacts in the final composite. Promising directions
    include a learned detection-confidence head tuned for this
    task, a track-lifetime filter that requires a detection to
    persist across several frames before entering the translation
    pipeline, and per-region saliency cues to suppress
    peripheral text.
  \item \textbf{Occlusion handling in compositing.} The current
    S5 compositing step paints the translated ROI over every pixel
    inside the target quad, including any foreground object---a
    hand, a passing vehicle, foliage---that happens to occlude the
    text surface in a given frame. A concrete extension reuses the
    Hi-SAM infrastructure we already vendor: segment foreground
    objects inside the target bbox, subtract their masks from the
    per-frame alpha mask so the composite never paints over them,
    and propagate the mask across frames with the same CoTracker
    machinery used for quad corners.
  \item \textbf{Temporal smoothing of the refiner's predictions.}
    The Alignment Refiner currently predicts $\Delta H$ per frame
    independently, leaving small inter-frame micro-oscillations.
    Smoothing over a short temporal window of predicted corner
    offsets is a natural extension.
  \item \textbf{Streaming main pipeline for long videos.} The main
    pipeline currently loads all video frames into memory, which
    limits it to short clips. A streaming architecture already
    exists in our data-generation tool; porting it to the main
    pipeline would lift this limit and extend the system to
    long-form content.
\end{itemize}
